%============================================================
%  app_algorithms.tex
%  Appendix D: Computational Aspects and Algorithms
%
%  Tensor Formats in Banach Spaces
%  Antonio Falcó, CEU Universities
%============================================================

\chapter{Computational Aspects and Algorithms}
\label{app:algorithms}

The theory developed in this monograph --- minimal subspaces, Tucker
and tree-based manifolds, the Dirac--Frenkel variational principle ---
has direct computational consequences. This appendix describes the
main algorithms used in practice, explains their connection to the
mathematical structures of the main text, and identifies precisely
where the Banach-space generalisations of
Chapters~\ref{chap:minsubspaces}--\ref{chap:DF} modify or constrain
the algorithmic picture.

Each algorithm is stated with pseudocode, its cost is analysed, and
its relationship to the abstract theory is made explicit. The principal
references are~\cite{Hackbusch2019} (systematic treatment of tensor
numerics), \cite{Bachmayr2016} (survey of tensor network methods), and
\cite{UschmajewVandereycken2020} (Riemannian optimisation on tensor
manifolds).

%------------------------------------------------------------
\section{Storage complexity of tensor formats}
\label{app:alg:storage}
%------------------------------------------------------------

\subsection*{The curse of dimensionality, quantified}

Let $V_j = \RR^n$ for $j = 1, \ldots, d$. The full tensor space
$\VD = \bigotimes_{j=1}^d \RR^n$ has dimension $n^d$: intractable
for moderate $n$ and $d$. Tensor formats replace this by structured
low-dimensional parametrisations. The table below summarises storage
cost in terms of maximal rank $r$ and mode size $n$:

\medskip
\begin{center}
\renewcommand{\arraystretch}{1.4}
\begin{tabular}{@{}llll@{}}
\toprule
\textbf{Format} & \textbf{Parameters}
    & \textbf{Storage} & \textbf{Tree}\\
\midrule
Full & $n^d$ entries & $O(n^d)$ & ---\\
Tucker & core $r^d$ + $d$ factor matrices $n\!\times\! r$
    & $O(r^d + dnr)$ & $\TDtuck$\\
HT (balanced binary)
    & transfer tensors at $O(d)$ nodes, each $r\!\times\! r^2$
    & $O(dr^3 + dnr)$ & $\TDht$\\
TT (Tensor Train)
    & $d$ core matrices, each $r\!\times\! n\!\times\! r$
    & $O(dr^2n)$ & $\TDtt$\\
\bottomrule
\end{tabular}
\end{center}
\medskip

\begin{remark}[Exponential vs.\ polynomial scaling]
\label{rem:app:scaling}
The Tucker format reduces factor storage to $O(dnr)$ (linear in $d$)
but retains the core at $O(r^d)$. The HT and TT formats remove this
remaining exponential cost by imposing rank constraints at each level
of the tree, reducing core storage to $O(dr^3)$ and $O(dr^2n)$
respectively. This is precisely the content of the intersection
representation (Theorem~\ref{thm:intersection_representation}):
$\FTv$ imposes rank constraints at \emph{every} level of $\TD$,
not only at the leaves.
\end{remark}

\subsection*{Parameter count for tree-based formats}

For a general dimension partition tree $\TD$ and rank tuple $\fr$,
the total parameter count is:
\begin{equation}
\label{eq:app:param_count}
    P(\fr,\TD)
    = \sum_{j \in D} r_{\{j\}} n_j
    + \sum_{\substack{\alpha \in \TD \setminus (\Leaves \cup \{D\})}}
      r_\alpha \prod_{\beta \in \Sons{\alpha}} r_\beta
    + \prod_{\beta \in \Sons{D}} r_\beta,
\end{equation}
where the three terms count, respectively, the leaf factor matrices,
the internal transfer tensors, and the root tensor.

%------------------------------------------------------------
\section{The Higher-Order SVD}
\label{app:alg:hosvd}
%------------------------------------------------------------

\subsection*{Setup and goal}

The \emph{Higher-Order SVD} (HOSVD) of De Lathauwer, De Moor, and
Vandewalle~\cite{DeLathauwer2000} computes a Tucker approximation.
Given $\bv \in \VD = \bigotimes_{j=1}^d \RR^{n_j}$ and target ranks
$\fr = (r_j)_{j \in D}$, it produces:
\begin{enumerate}
\item[(i)] Factor matrices $U_j \in \RR^{n_j \times r_j}$ with
    orthonormal columns, approximating $\Umin{j}{\bv}$.
\item[(ii)] A core tensor
    $C^{(D)} \in \RR^{r_1 \times \cdots \times r_d}$ such that
    $\bv \approx \sum_{(i_j)} C^{(D)}_{(i_j)}
    \bigotimes_j u_{i_j}^{(j)}$.
\end{enumerate}

\subsection*{Algorithm}

\noindent\textbf{Algorithm D.1} (HOSVD).\\
\textit{Input:} tensor $\bv \in \RR^{n_1 \times \cdots \times n_d}$,
target ranks $(r_1,\ldots,r_d)$.\\
\textit{Output:} core tensor $C^{(D)}$, factor matrices
$U_1,\ldots,U_d$.
\begin{enumerate}
\item \textbf{Mode unfoldings.} For each $j = 1,\ldots,d$:
    \begin{enumerate}
    \item Compute the $j$-mode unfolding
        $M_j = \Mmat{\{j\}}(\bv) \in \RR^{n_j \times \prod_{k \neq j} n_k}$.
    \item Compute the truncated SVD: $M_j \approx U_j \Sigma_j V_j^\top$,
        retaining the $r_j$ largest singular values.
        Set $U_j \in \RR^{n_j \times r_j}$ (left singular vectors).
    \end{enumerate}
\item \textbf{Core compression.} Compute
    \begin{equation*}
        C^{(D)}_{i_1,\ldots,i_d}
        = \sum_{k_1,\ldots,k_d}
          (U_1)_{k_1 i_1} \cdots (U_d)_{k_d i_d}\;
          \bv_{k_1,\ldots,k_d},
    \end{equation*}
    i.e., project $\bv$ onto the subspaces $\spann(U_j)$ in each mode.
\item \textbf{Output.} Return $(C^{(D)}, U_1,\ldots,U_d)$.
\end{enumerate}

\subsection*{Cost and approximation quality}

The dominant cost is $d$ SVDs of unfolding matrices of size
$n_j \times \prod_{k \neq j} n_k$, giving arithmetic cost
$O(dn^{d+1})$, which still scales exponentially in $d$.
For tensors already in a low-rank format, recompression costs
$O(dr^{d+1})$ or $O(dr^3)$ depending on the tree structure.

\begin{theorem}[HOSVD quasi-optimality; {\cite[Theorem~4.2]{DeLathauwer2000}}]
\label{thm:app:hosvd}
Let $\bv \in \VD$ and let $\bw$ be the output of \textup{Algorithm~D.1}
with ranks $(r_j)_j$. Then
\begin{equation}
\label{eq:app:hosvd_bound}
    \|\bv - \bw\|_F
    \leq \sqrt{d}\;
    \min_{\bu \in \Tucker{\leq\fr}{\VD}} \|\bv - \bu\|_F.
\end{equation}
The HOSVD is not a best-approximation algorithm, but provides a
quasi-optimal approximation with factor $\sqrt{d}$.
\end{theorem}

\begin{remark}[Connection to minimal subspaces]
\label{rem:app:hosvd_min}
The factor matrix $U_j$ is the best rank-$r_j$ approximation to
the column space of the $j$-mode unfolding $M_j$, whose column space
is exactly $\Umin{j}{\bv}$ when $r_j = r_j(\bv)$ (the exact rank).
Algorithm~D.1 is therefore a computational procedure for approximating
the minimal subspaces of Chapter~\ref{chap:minsubspaces}.

In the Banach setting (non-Hilbert $V_j$), the SVD is replaced by a
best rank-$r_j$ approximation in the operator norm of
$\cL(V_j^*,\Vbigalpha^{[j]})$. Its existence is guaranteed by
Theorem~\ref{thm:best_approx_tucker}, but the quasi-optimality
factor $\sqrt{d}$ may worsen depending on the Banach--Mazur distance
of $\VDnorm$ from a Hilbert space.
\end{remark}

%------------------------------------------------------------
\section{Alternating Least Squares}
\label{app:alg:als}
%------------------------------------------------------------

\subsection*{The ALS problem and microiterations}

Given $\bu \in \VDnorm$ and rank tuple $\fr$, the best Tucker
approximation problem is
\begin{equation}
\label{eq:app:als_problem}
    \min_{\bw \in \Tucker{\leq\fr}{\VD}} \|\bu - \bw\|_D^2.
\end{equation}
\emph{Alternating Least Squares} (ALS) solves this by fixing all
factor matrices but one and optimising over the remaining one,
cycling over all modes.

\medskip
\noindent\textbf{Algorithm D.2} (ALS for Tucker format).\\
\textit{Input:} $\bu \in \VDnorm$, ranks $\fr$, initial factor
matrices $\{U_j^{(0)}\}_{j=1}^d$.\\
\textit{Repeat until convergence:}
\begin{enumerate}
\item For $j = 1,\ldots,d$:
    \begin{enumerate}
    \item \textbf{Fix} all $U_k$ for $k \neq j$ at their current values.
    \item \textbf{Form the reduced right-hand side:} compute the array
        \begin{equation*}
            \bigl(R_j\bigr)_{i_j, (i_k)_{k \neq j}}
            \coloneqq
            \sum_{(k_l)} \bu_{k_1,\ldots,k_d}
            \prod_{l \neq j} (U_l)_{k_l, i_l},
        \end{equation*}
        which is the tensor $\bu$ projected onto all modes except $j$.
        This gives $R_j \in \RR^{n_j \times \prod_{k \neq j} r_k}$.
    \item \textbf{Solve the mode-$j$ subproblem:}
        \begin{equation*}
            U_j^{\mathrm{new}}
            = \argmin_{U \in \RR^{n_j \times r_j}}
              \bigl\|R_j - U\, H_j\bigr\|_F^2,
        \end{equation*}
        where $H_j \in \RR^{r_j \times \prod_{k \neq j} r_k}$ is the
        Gram factor assembled from the other modes.
    \item \textbf{Orthonormalise:} $U_j \leftarrow U_j^{\mathrm{new}}$
        after a QR step.
    \end{enumerate}
\item \textbf{Update the core:}
    $C^{(D)}_{i_1,\ldots,i_d} \leftarrow
    \sum_{k_1,\ldots,k_d} \bu_{k_1,\ldots,k_d}
    \prod_j (U_j)_{k_j,i_j}$.
\end{enumerate}

\subsection*{Convergence}

\begin{theorem}[ALS convergence in Hilbert spaces;
    {\cite[Theorem~4.1]{UschmajewVandereycken2020}}]
\label{thm:app:als_convergence}
Let $\VDnorm$ be a Hilbert tensor space and let $(\bw^{(k)})$ be
the sequence generated by \textup{Algorithm~D.2}. Then:
\begin{enumerate}
\item[\textup{(i)}] The residuals $\|\bu - \bw^{(k)}\|$ are monotonically
    non-increasing.
\item[\textup{(ii)}] Every accumulation point $\bw^*$ satisfies
    the first-order optimality conditions in each mode direction:
    $\langle \bu - \bw^*, v \otimes \bw^*_{[j]}\rangle = 0$
    for all $v \in V_j$, $j = 1,\ldots,d$.
\item[\textup{(iii)}] Global convergence to a best approximation is
    not guaranteed; ALS may converge to a local minimum or saddle point.
\end{enumerate}
\end{theorem}

\begin{remark}[ALS in Banach spaces]
\label{rem:app:als_banach}
In a reflexive Banach tensor space $\VDnorm$, the mode-$j$ subproblem
in step~1(c) becomes: find the best rank-$r_j$ operator approximation
in the norm $\|\cdot\|_D$ projected onto mode~$j$. This is a convex
problem whose unique solution exists by
Theorem~\ref{thm:best_approx_tucker}. Monotone decrease of part~(i)
holds in full generality. The characterisation of accumulation points
in part~(ii) must be rephrased using the duality mapping $J$
(Appendix~\ref{app:functional}) in place of the inner product:
\begin{equation*}
    \langle v \otimes \bw^*_{[j]},\,
    J(\bu - \bw^*)\rangle = 0
    \quad \forall\, v \in V_j,\; j = 1,\ldots,d.
\end{equation*}
The convergence analysis in this Banach setting is Open
Problem~\ref{prob:ALS}.
\end{remark}

%------------------------------------------------------------
\section{The projector-splitting integrator}
\label{app:alg:splitting}
%------------------------------------------------------------

\subsection*{The Dirac--Frenkel ODE}

Consider the Dirac--Frenkel reduced model on $\Tucker{\fr}{\VD}$
(Chapter~\ref{chap:DF}):
\begin{equation}
\label{eq:app:df_reduced}
    \dot{\bv}_r(t)
    = \Pmetric{\bv_r(t)}\bigl(\bF(t,\bv_r(t))\bigr),
    \quad \bv_r(0) = \bv_0.
\end{equation}
Direct integration requires evaluating $\Pmetric{\bv_r(t)}$ at each
step. The \emph{projector-splitting integrator}
of~\cite{LubichOseledetsVandereycken2015} avoids this by decomposing
the projection into a sequence of simpler updates.

\subsection*{The TT case}

For the TT format ($\TD = \TDtt$, bond dimensions
$r_1,\ldots,r_{d-1}$), the tangent space projection splits as
\begin{equation}
\label{eq:app:splitting}
    \Pmetric{\bv_r}
    = \sum_{k=1}^{d} P_k
      - \sum_{k=1}^{d-1} P_{k,k+1},
\end{equation}
where $P_k$ projects onto the subspace with the $k$-th component free
and $P_{k,k+1}$ is the correction at the $(k,k{+}1)$ interface.

\medskip
\noindent\textbf{Algorithm D.3} (Projector-splitting, one step $h$).\\
\textit{Input:} $\bv_r^n$ in TT format, vector field $\bF$, step $h$.\\
\textit{Output:} $\bv_r^{n+1}$ in TT format.
\begin{enumerate}
\item \textbf{Forward sweep} ($k = 1,\ldots,d$): integrate the
    $k$-th component ODE
    \begin{equation*}
        \dot{C}^{(k)}(t)
        = \bigl\langle \bF(t,\bv_r(t)),\,
          P_k(\,\cdot\,) \bigr\rangle
    \end{equation*}
    from $t^n$ to $t^n + h$; left-orthogonalise via QR.
\item \textbf{Correction sweep} ($k = d-1,\ldots,1$): apply the
    correction steps $-P_{k,k+1}$ symmetrically.
\item \textbf{Return} the updated TT tensor $\bv_r^{n+1}$.
\end{enumerate}

\begin{theorem}[Error bound in Hilbert spaces;
    {\cite[Theorem~4.1]{LubichOseledetsVandereycken2015}}]
\label{thm:app:splitting_error}
Let $\VDnorm$ be a Hilbert tensor space and $\bF$ Lipschitz with
constant $L$. Then there exists $C > 0$ (depending on $L$, the
curvature of $\Tucker{\fr}{\VD}$, and the initial data) such that
\begin{equation}
\label{eq:app:splitting_bound}
    \|\bu(t^n) - \bv_r^n\|
    \leq Ch + O\!\bigl(\delta_{\fr}(\bu)\bigr),
\end{equation}
where $\delta_{\fr}(\bu)
= \inf_{\bw \in \Tucker{\leq\fr}{\VD}} \|\bu - \bw\|$ is the
best-approximation error.
\end{theorem}

The curvature term in $C$ is the quantity whose Banach-space
analogue appears in Open Problem~\ref{prob:integrator}. For general
tree-based formats, the splitting integrator
of~\cite{CerutiLubichWalach2020} applies a level-by-level sweep along
$\TD$ at cost $O(\dep{\TD} \cdot dr^3)$ per time step.

%------------------------------------------------------------
\section{DMRG and the Tucker manifold}
\label{app:alg:dmrg}
%------------------------------------------------------------

\subsection*{DMRG as ALS on the TT manifold}

The \emph{Density Matrix Renormalisation Group} (DMRG), introduced by
White~\cite{White1992}, is the most widely used algorithm for
ground-state computation in quantum many-body systems. In its modern
formulation~\cite{HoltzRohwedderSchneider2012}, DMRG is equivalent to
ALS on the TT manifold applied to the eigenvalue problem
$H\bu = \lambda\bu$: each microiteration updates one TT component
while fixing all others, producing an effective eigenvalue problem
of size $r^2 n$.

\begin{remark}[Gauge conditions and the Tucker bundle]
\label{rem:app:dmrg}
The \emph{left-orthogonal gauge} in DMRG requires each component tensor
$C^{(k)} \in \RR^{r_{k-1} \times n_k \times r_k}$ to satisfy
\begin{equation*}
    \sum_{i_k=1}^{n_k}
    \bigl(C^{(k)}\bigr)_{:,i_k,:}^{\!\top}
    \bigl(C^{(k)}\bigr)_{:,i_k,:}
    = I_{r_k}.
\end{equation*}
This is the computational analogue of fixing the complement $W_k$ at
each tree node to be the orthogonal complement of $\spann(C^{(k)})$,
giving a canonical trivialisation of the Tucker bundle
(Section~\ref{sec:tucker_manifold:bundle}).
The gauge group $\prod_k \GL{r_k}{\RR}$ of
Remark~\ref{rem:gauge_freedom} is quotiented out by the QR
orthogonalisation step.
\end{remark}

\subsection*{Two-site DMRG and rank adaptation}

One-site DMRG cannot change the bond dimensions $r_k$. \emph{Two-site
DMRG} optimises two adjacent components simultaneously, forming an
effective tensor of rank $r_k \cdot r_{k+1}$ which is then truncated
by SVD to obtain updated bond dimensions. Geometrically, this
corresponds to moving between strata of the stratification
$\VD = \bigsqcup_\fr \Tucker{\fr}{\VD}$
(Theorem~\ref{thm:tucker_decomposition_full}).

%------------------------------------------------------------
\section{Retraction-based optimisation on tensor manifolds}
\label{app:alg:retraction}
%------------------------------------------------------------

\subsection*{Riemannian gradient methods}

In the Hilbert case, the Tucker manifold is a Riemannian submanifold
of $\VDnorm$. Riemannian gradient descent~\cite{UschmajewVandereycken2020}
for minimising $f: \VDnorm \to \RR$ over $\Tucker{\fr}{\VD}$ reads:
\begin{equation}
\label{eq:app:riem_gradient}
    \bv_{k+1}
    = \mathrm{Ret}_{\bv_k}\!\bigl(
        -\alpha_k\,\Pmetric{\bv_k}(\nabla f(\bv_k))
    \bigr),
\end{equation}
where $\Pmetric{\bv}(\nabla f(\bv))$ is the Riemannian gradient and
$\mathrm{Ret}_\bv: \ZD{\bv} \to \Tucker{\fr}{\VD}$ is a retraction.

\begin{definition}[Retraction]
\label{def:app:retraction}
A \emph{retraction} on a Banach manifold $\mathbb{M}$ is a
$\mathcal{C}^1$-map $\mathrm{Ret}: T\mathbb{M} \to \mathbb{M}$
such that, for each $m \in \mathbb{M}$:
\begin{enumerate}
\item[\textup{(i)}] $\mathrm{Ret}_m(0_m) = m$,
\item[\textup{(ii)}] $D(\mathrm{Ret}_m)(0_m) = \Id{T_m\mathbb{M}}$.
\end{enumerate}
\end{definition}

\begin{example}[Tucker retraction via HOSVD truncation]
\label{ex:app:tucker_retraction}
A practical retraction on $\Tucker{\fr}{\VD}$ is:
\begin{equation*}
    \mathrm{Ret}_{\bv}(\dot\bv)
    = \mathrm{HOSVD}_\fr(\bv + \dot\bv),
\end{equation*}
i.e., apply Algorithm~D.1 to $\bv + \dot\bv$ and truncate to rank
$\fr$. This satisfies (i) and (ii) with quasi-optimality constant
$\sqrt{d}$ (Theorem~\ref{thm:app:hosvd}), and remains valid in Banach
spaces as long as the best-approximation step is solvable ---
guaranteed by Theorem~\ref{thm:best_approx_tucker}.
\end{example}

\subsection*{The Banach case: duality-mapping gradient}

In a non-Hilbert Banach tensor space $\VDnorm$, the Riemannian gradient
is replaced by the \emph{duality-mapping gradient}: the unique element
$\nabla_B f(\bv) \in \ZD{\bv}$ satisfying
\begin{equation}
\label{eq:app:banach_gradient}
    \bigl\langle \nabla_B f(\bv),\,
    J\!\bigl(\bv - \bu\bigr)\bigr\rangle
    = Df(\bv)[\bv - \bu]
    \qquad \forall\,\bu \in \Tucker{\fr}{\VD},
\end{equation}
where $J: \VDnorm \to (\VDnorm)^*$ is the normalised duality mapping
(Appendix~\ref{app:functional}, \S\ref{app:functional:index}).
The iteration~\eqref{eq:app:riem_gradient} becomes a generalised
gradient step using the generalised projection $\Pi_{\bv}$ of
Appendix~\ref{app:functional} in place of $\Pmetric{\bv}$. Convergence
analysis in this setting is related to Open Problem~\ref{prob:finsler}.

%------------------------------------------------------------
\section{Historical notes}
\label{app:alg:history}
%------------------------------------------------------------

\paragraph{Tucker decomposition and HOSVD.}
The Tucker decomposition~\cite{Tucker1966} was introduced as a
psychometric tool and has been used in data analysis since the 1970s.
The HOSVD as an algorithm for best Tucker approximations was introduced
by De Lathauwer, De Moor, and Vandewalle~\cite{DeLathauwer2000}; the
quasi-optimality bound~\eqref{eq:app:hosvd_bound} is from the same
paper. A comprehensive survey of tensor decompositions is given in
Kolda and Bader~\cite{KoldaBader2009}.

\paragraph{Tensor Train and DMRG.}
The TT decomposition was introduced independently by Vidal~\cite{Vidal2003}
(matrix product states, MPS) and by Oseledets and
Tyrtyshnikov~\cite{Oseledets2011}. The DMRG algorithm of
White~\cite{White1992} predates both and was recognised
in~\cite{HoltzRohwedderSchneider2012} as ALS on the TT manifold.
The TT-SVD of~\cite{Oseledets2011} provides quasi-optimal TT
approximations analogous to the HOSVD for Tucker.

\paragraph{Projector-splitting integrators.}
The projector-splitting integrator for TT formats was introduced by
Lubich, Oseledets, and Vandereycken~\cite{LubichOseledetsVandereycken2015},
building on the matrix case of Koch and Lubich~\cite{KochLubich2007}
(MCTDH). The extension to general tree-based formats is
in~\cite{CerutiLubichWalach2020}. Convergence in the Banach setting
is Open Problem~\ref{prob:integrator}.

\paragraph{Riemannian optimisation.}
Riemannian gradient methods on Tucker and TT manifolds were
systematically developed by Uschmajew and
Vandereycken~\cite{UschmajewVandereycken2020}, using the Riemannian
structure in a way that parallels
Chapters~\ref{chap:tucker_manifold}--\ref{chap:TB_manifold} but
restricted to the Hilbert case. The Banach generalisation is related
to Open Problem~\ref{prob:finsler}.
